\chapter{Derivatives and the Binder Problem}\label{ch:deriv}

The derivative rules were the first rules that could not simply be added to the existing story.
Their termination broke the additive measure (Part~III), and their \emph{meaning} broke the
semantics of the previous chapter in a way that is worth understanding, because it is the
reason the engine has three semantics rather than one.

\section{The rules}

$\texttt{diff}(f, x)$ is a function node, \lc{.fn "diff" [f, .var x]}, and the derivative rules
push it inward: the derivative of a sum is the sum of the derivatives, of a product the product
rule, of a power the power rule, of $\sin u$ the chain rule with $\cos u$ outside. Each is one
rewrite, each is recorded, and the trace of $\frac{d}{dx}(x^2 \sin x)$ reads like a textbook
because it \emph{is} the textbook's derivation, one rule per line.

\section{Why \lc{evalR} could not take the case}

The obvious plan was to extend \lc{evalR} with a case for \texttt{diff}. It cannot be done, and
the reason is the second child.

In $\texttt{diff}(f, x)$ the argument $x$ is not a value; it is a \emph{binder}, the name of the
variable we differentiate along. But \lc{Expr} has no binder positions --- every child of every
node is a term. So \lc{.var x} and \lc{.add [.var x]} denote the same real number, while only one
of them can be read as ``the variable we differentiate along''. A semantics that interprets
\texttt{diff} therefore cannot satisfy the \lc{congr} law of Chapter~\ref{ch:sound}: rewriting
the child \lc{.add [.var x]} to \lc{.var x} is sound for every other node and changes the meaning
of this one. And \lc{congr} is the law the fold rests on.

\section{A third semantics}

So the semantics was extended rather than edited. \lc{evalD} agrees with \lc{evalR} everywhere
except that it reads \texttt{diff}, via Mathlib's \lc{deriv}:

\begin{minted}{lean}
def fnD : EnvR → String → List Expr → ℝ
  | ρ, f, [a] => applyFn f (evalD ρ a)
  | ρ, f, [g, v] =>
    if f = "diff" then
      match v with
      | .var x => deriv (fun t => evalD (upd ρ x t) g) (ρ x)
      | _ => 0
    else 0
  | _, _, _ => 0
\end{minted}

The theorem \lc{evalD_eq_evalR} says the two agree on every term without a \texttt{diff} node.
This is the third layer of one pattern: \lc{eval?} $\subset$ \lc{evalR} $\subset$ \lc{evalD},
each proved a restriction of the next. And because the derivative rules never run under the
verified \lc{normalize} of Chapter~\ref{ch:rewrite} --- they run under the tiered rewriter of
Part~III --- they never needed the congruence law in the first place. Each rule gets its own
theorem, at the whole-term level, with \lc{deriv}.

\section{Which rules need a hypothesis}

Here the junk-value discipline of the previous chapter earns its keep. \lc{deriv} is junk-valued
where a function is not differentiable, so the sum, product and chain rules genuinely fail
without a hypothesis, and the hypothesis has to be stated.

\begin{itemize}[nosep]
  \item \status{verified}: \rn{diff.constant}, \rn{diff.variable}, \rn{diff.constant-multiple}.
    Pulling a constant factor out needs no differentiability at all.
    The Mathlib lemma is \lc{deriv_const_mul_field}.
  \item \status{conditional}: \rn{diff.sum} (every summand differentiable at the point),
    \rn{diff.product} (both factors), \rn{diff.power} (a natural exponent and a differentiable
    base; the real-exponent case is open), \rn{diff.chain} (the inner function differentiable;
    proved for $\sin$, $\cos$, $\exp$; $\ln$ and $\tan$ also need a domain condition and are open).
\end{itemize}

And the necessity is proved, as it was for the algebraic rules:

\begin{thmbox}[\lc{not_diff_sum_sound}]
  The sum rule is not unconditional: at $0$, the rule gives $|x| + x$ the derivative $1$, where
  the derivative does not exist.
\end{thmbox}

The engine reports all of this through the capabilities message, so the notebook's panel says
``conditional: needs both factors differentiable at the point'' instead of ``no soundness
theorem yet''. That message was the first place where the statuses of the Preface became
visible to a user.

\begin{logbox}[Not claimed]
  \rn{diff.higher-order} is an abbreviation --- it eliminates the three-argument form
  $\texttt{diff}(f, x, n)$ into $n$ nested derivatives --- so there is nothing to prove.
  \rn{diff.matrix} is proved but vacuous: a matrix has no value in the real semantics, so the
  theorem has no content. The honest label for both is \status{unverified}, and the note next to
  the label says why.
\end{logbox}

\begin{trybox}
\begin{minted}{text}
diff(abs(x) + x, x)
diff(x^x, x)
diff(ln(x), x)
\end{minted}
  The first is the counterexample of \lc{not_diff_sum_sound} in the notebook: the sum rule
  fires, leaves the derivative of $|x|$ standing, and its status says what it assumed. The second uses a rule the theorems do not yet cover
  (a real exponent) and says so.
\end{trybox}

\section{Exercises}

\begin{exercisebox}[The binder problem, concretely]
  Give two terms that \lc{evalR} cannot distinguish but that a semantics for \texttt{diff} must,
  and show which law of \lc{Congruence} they break.
\end{exercisebox}

\begin{exercisebox}[A missing case]
  The chain rule for $\ln u$ needs $u > 0$ at the point (and differentiability). Write the
  theorem statement in the style of \lc{diff_chain_sin_sound}, and say which Mathlib lemma
  about \lc{Real.log} you would need.
\end{exercisebox}
